% ElegantNote — a lightweight lecture/reading-notes class (article-based,
% not chapter-structured), companion to ElegantBook. Ships with TeX Live;
% not bundled.
%
% backend=bibtex is passed explicitly: the class defaults to biber, which
% this catalogue's compile loop does not run (only classic bibtex).
%
% Class documentation: texdoc elegantnote | https://ctan.org/pkg/elegantnote
% Licence: LPPL-1.3c or later, (c) Dongsheng Deng, Ran Wang.
\documentclass[en,blue,a4,10pt,bibtex]{elegantnote}

\title{Notes on Ordinary Differential Equations}
\author{Priya Anand}
\institute{Department of Applied Mathematics}
\version{1.0}
\date{\today}

\addbibresource[location=local]{reference.bib}

\begin{document}
\maketitle

\section{First-Order Linear Equations}

A first-order linear ODE has the form $y' + p(x)y = q(x)$. Multiplying
by the integrating factor $\mu(x) = e^{\int p(x)\,dx}$ turns the left
side into an exact derivative:
\[
\frac{d}{dx}\big[\mu(x) y\big] = \mu(x) q(x).
\]

\begin{remark}
This is the single most useful trick for first-order linear equations
— everything else in this section is a variation on it.
\end{remark}

\begin{example}
Solve $y' + 2xy = x$. Here $\mu(x) = e^{x^2}$, so
$\big(e^{x^2} y\big)' = x e^{x^2}$, giving
$e^{x^2} y = \tfrac{1}{2} e^{x^2} + C$, i.e.
$y = \tfrac{1}{2} + C e^{-x^2}$.
\end{example}

\section{Second-Order Constant-Coefficient Equations}

For $ay'' + by' + cy = 0$, the characteristic equation
$a\lambda^2 + b\lambda + c = 0$ determines the solution's shape:

\begin{itemize}
  \item two real roots $\lambda_1 \neq \lambda_2$: $y = C_1 e^{\lambda_1 x} + C_2 e^{\lambda_2 x}$;
  \item a repeated root $\lambda$: $y = (C_1 + C_2 x) e^{\lambda x}$;
  \item complex roots $\alpha \pm i\beta$: $y = e^{\alpha x}(C_1 \cos \beta x + C_2 \sin \beta x)$.
\end{itemize}

\begin{note}
The repeated-root case is exactly the boundary between oscillation and
pure exponential decay/growth — it is worth sketching all three cases
on the same axes to see why.
\end{note}

\section{Numerical Methods}

When no closed form exists, Euler's method approximates the solution
of $y' = f(x, y)$, $y(x_0) = y_0$ by stepping
$y_{n+1} = y_n + h f(x_n, y_n)$. It has local truncation error
$O(h^2)$ and global error $O(h)$, per \textcite{butcher2016numerical}.
The classical fourth-order Runge--Kutta method improves this to global
error $O(h^4)$ at the cost of four function evaluations per step,
following \textcite{hairer1993solving}.

\printbibliography

\end{document}
